% ============================================================================
%  5. DATOS Y CIERRE  (165--180 min)
% ============================================================================
\section{Cierre}

% S51
\begin{frame}[label=vs-datos]{El modelo frente a los datos}
  \centering\small
  \renewcommand{\arraystretch}{1.25}
  \begin{tabular}{@{}lcccc@{}}
    \toprule
    & \multicolumn{2}{c}{\textbf{Vol.\ relativa a $y$}} & \multicolumn{2}{c}{\textbf{Corr.\ con $y$}}\\
    \cmidrule(lr){2-3}\cmidrule(lr){4-5}
    & Modelo & Datos & Modelo & Datos\\
    \midrule
    Consumo & 0.44 & 0.53 & 0.94 & 0.82\\
    Inversión & 3.38 & 3.91 & 0.99 & 0.93\\
    Horas & \num{0.49} & \num{1.23} & 0.97 & 0.87\\
    Salario & 0.54 & 0.63 & \num{0.98} & \num{0.22}\\
    \bottomrule
  \end{tabular}
  \par\vspace{0.3cm}\raggedright\normalsize
  Componente cíclico, filtro HP con $\lambda=1600$: en Dynare basta añadir
  \texttt{hp\_filter = 1600} a \texttt{stoch\_simul}.\par
  \vspace{0.2cm}
  Acierta el orden de volatilidades: la inversión es mucho más volátil que el producto
  y el consumo, más suave.
  \fuente{Filtro de \textcite{hp1997}. Datos: EE.UU.\ 1954--2015, \textcite[Tabla~1.2]{gmo}. Modelo: momentos teóricos
  de Dynare. Handout, Cuadro~3.15.}
  \note{$\sigma_y$ del modelo filtrado: 1.4\%, frente a 1.5\% en los datos.}
\end{frame}

% S52
\begin{frame}[label=limites]{Dos cosas que este RBC no logra explicar}
  \begin{columns}[T,onlytextwidth]
    \begin{column}{0.48\textwidth}
      \textbf{1. Poca propagación propia}\par
      \vspace{0.2cm}
      $\hat y/\hat a$: $1.48$ en $t=1$; $1.74$ en $t=40$.\par
      \vspace{0.2cm}
      La forma de $\hat y$ copia la de $\hat a$: casi toda la persistencia viene de
      $\phi=0.979$, no del modelo.
      \fuente{\textcite{cogleynason1995}.}
    \end{column}
    \begin{column}{0.48\textwidth}
      \textbf{2. Horas poco volátiles}\par
      \vspace{0.2cm}
      Modelo: $0.49$ veces el producto. Datos: $1.23$.\par
      \vspace{0.2cm}
      Aun con sustitución intertemporal, las horas se mueven poco. Motiva el trabajo
      indivisible y otras extensiones.
      \fuente{\textcite{hansen1985}.}
    \end{column}
  \end{columns}
  \par\vspace{0.5cm}
  ¿Y el salario? En el modelo es la productividad laboral, $w=\alpha y/h$: con un solo
  choque, todo se mueve al unísono.
  \note{Vuelve la pregunta de la slide impulso: la persistencia viene del choque, no de la
  propagación.}
\end{frame}

% S53
\begin{frame}[label=resumen]{La historia en una tabla}
  \centering\scriptsize
  \renewcommand{\arraystretch}{1.3}
  \setlength{\tabcolsep}{4pt}
  \begin{tabular}{@{}llll@{}}
    \toprule
    \textbf{Momento} & \textbf{Ecuación} & \textbf{Número} & \textbf{Mecanismo}\\
    \midrule
    Impacto: horas & $\hat h=1.715(\hat a-\hat c)$ & 0.515 & sustitución\\
    Impacto: producto & $\hat y=\hat a+\alpha\hat h$ & 1.064 & amplificación\\
    Impacto: reparto & $\hat y=0.798\hat c+0.202\hat\imath$ & 69\% invertido & suavizamiento\\
    Capital & $\Delta\hat k=\delta(\hat\imath-\hat k)$ & máx.\ en $t=26$ & acumulación lenta\\
    Consumo & $\Delta\hat c=\beta r_{ss}\hat r$ & máx.\ en $t=19$ & Euler\\
    Regreso & $\hat h<0$ desde $t=30$ & $-0.031$ en $t=40$ & efecto riqueza\\
    Contraste $\delta=1$ & $\hat y-\hat c=1.25\hat h$ & $\hat h=0$ & $c/y$ constante\\
    \bottomrule
  \end{tabular}
  \par\vspace{0.35cm}\raggedright\normalsize
  Todos los números: la receta o \texttt{oo\_.irfs} en Dynare.
\end{frame}

% S54
\begin{frame}[fragile,label=estimacion]{Siguiente paso: estimar en vez de calibrar}
  Hoy fijamos los parámetros. Dynare también puede \textbf{estimarlos} con datos:\par
  \vspace{0.25cm}
  \begin{enumerate}
    \item La solución se escribe en forma estado-espacio.\par\vspace{0.1cm}
    \item El filtro de Kalman da la verosimilitud de los datos.\par\vspace{0.1cm}
    \item Bayes: posterior $\propto$ verosimilitud $\times$ prior.\par\vspace{0.1cm}
    \item MCMC recorre la posterior.
  \end{enumerate}
  \vspace{0.2cm}
  \begin{lstlisting}
varobs y;
estimated_params; phi, beta_pdf, 0.9, 0.05; end;
estimation(datafile = datos, mh_replic = 20000);
\end{lstlisting}
  \fuente{Esquema ilustrativo. Handout, §3.14.}
  \note{Solo un adelanto: no correr en clase. La sintaxis completa y sus supuestos están en
  el handout.}
\end{frame}

% S55
\begin{frame}[label=material]{Para seguir practicando}
  \textbf{En el repositorio}\par
  \vspace{0.15cm}
  \small
  \begin{tabular}{@{}ll@{}}
    \texttt{handout/handout\_dynare.tex} & el handout completo\\
    \texttt{handout/rbc\_analytical.mod} & el modelo de referencia\\
    \texttt{handout/graficar\_irfs.m} & IRF con estilo\\
    \texttt{handout/reconstruir\_irf.m} & IRF a mano desde \texttt{oo\_.dr}\\
    \texttt{handout/cuatro\_formas/} & el mismo modelo, cuatro escrituras\\
    \texttt{taller/rbc\_taller.mod} & el taller de hoy
  \end{tabular}
  \par\vspace{0.4cm}\normalsize
  \textbf{Ejercicios del handout}\par
  \vspace{0.15cm}
  Barrido en $\phi$, segundo orden, filtro HP, romper el timing y las cuatro formas con
  $\delta=1$.
\end{frame}

% S56
\begin{frame}[label=referencias]{Referencias: método y herramientas}
  \renewcommand*{\bibfont}{\footnotesize}
  \printbibliography[heading=none,keyword=metodo]
\end{frame}

% S57
\begin{frame}[label=referencias-modelo]{Referencias: el modelo}
  \renewcommand*{\bibfont}{\footnotesize}
  \printbibliography[heading=none,keyword=modelo]
\end{frame}
